% TOP_PROBLEM: {"id": 2898, "title": "Kirby Problem 4.22", "category_id": 7, "category": {"id": 7, "name": "topology", "display_name": "Topology", "description": "Properties preserved under continuous deformations.", "slug": "topology", "order_index": 7, "created_at": "2026-07-31T15:26:25.671Z"}, "status": "open"}
% FIRST_POSTED: null
% ENTRY_KIND: statement_only
% Imported from: batch13.tex; UnsolvedMath ID 2898
\documentclass[11pt]{article}
\usepackage[margin=25mm]{geometry}
\usepackage{fontspec}
\setmainfont{Latin Modern Roman}
\usepackage{amsmath,amssymb,mathtools}
\usepackage{unicode-math}
\setmathfont{Latin Modern Math}
\usepackage{xurl,hyperref}
\hypersetup{colorlinks=true,urlcolor=blue}
\setcounter{secnumdepth}{0}
\setlength{\parindent}{0pt}
\setlength{\parskip}{5pt}
\setlength{\emergencystretch}{3em}
\newcommand{\sref}[2]{\hyperlink{source:#1:#2}{[S#2]}}
\newcommand{\cataloguescope}[1]{\paragraph{Recorded target.}#1}
\begin{document}
% UnsolvedMath ID 2898; source catalogue code KP-4.22
\setcounter{section}{2897}
\section{Kirby Problem 4.22}\label{2898}
{\small\sffamily UnsolvedMath ID 2898 · Topology}
\subsection{Definitions and mathematical statement}

\paragraph{Shared notation.}$\mathbb N=\{1,2,\ldots\}$, and $\mathbb Z,\mathbb Q,\mathbb R,\mathbb C$ denote the integers, rationals, reals and complex numbers. Any other conventions are specified locally.


Let $X$ be a closed smooth four-manifold with $\pi_1(X)=\pi$ a good group; orientability is not assumed.

A group $\pi$ is good in the Freedman--Quinn disk-embedding sense. The accepted capped-grope criterion is: for every based height-$1.5$ disk-like capped grope $G$ and homomorphism $\rho:\pi_1(G)\to\pi$, there is an immersed disk in $G$ with the same framed attaching boundary whose double-point loops all map to $1$ under $\rho$. A capped grope is the framed four-dimensional thickening of a surface complex built by attaching higher-stage surfaces along symplectic-basis curves and capping its tips by disks; height $1.5$ means a first stage and a second stage on one half of a symplectic basis. A double-point loop follows the two preimage branches of a self-intersection, with paths to the basepoint. The source gives this criterion in Problem 4.46, Remarks (1)--(2).

For a closed four-manifold $X$, a marked structure is a pair $(Y,f)$ with $f:Y\to X$ a homotopy equivalence carrying the local fundamental class to that of $X$. In the smooth structure set $\mathcal S_{\mathrm{DIFF}}(X)$, two pairs are equivalent when a diffeomorphism of their source manifolds makes the maps to $X$ homotopic; in $\mathcal S_{\mathrm{TOP}}(X)$ use a homeomorphism instead. The forgetful map $F$ keeps the underlying topological marked structure. Let $w:\pi_1(X)\to\{1,-1\}$ be the orientation character. The involution on $\mathbb Z[\pi_1(X)]$ is $\overline g=w(g)g^{-1}$, extended linearly. Write $L_5^s(\mathbb Z[\pi_1(X)],w)$ for Wall's degree-five simple surgery-obstruction group, with this involution; the decoration $s$ records simple, rather than merely homotopy, surgery data. The Wall realization action constructs topological marked structures by a five-dimensional normal cobordism with the specified surgery obstruction.

For $n\ge2$ and $\gcd(n,q)=1$, let $L(n,q)$ be the lens space obtained by quotienting $S^3\subset\mathbb C^2$ by $(z_1,z_2)\mapsto(\zeta z_1,\zeta^qz_2)$, where $\zeta=e^{2\pi i/n}$. A fake $S^1\times L(n,q)$ is homotopy equivalent but not homeomorphic to that product. In the dihedral case, $\operatorname{UNil}_1(\mathbb Z;\mathbb Z^-,\mathbb Z^-)$ denotes Cappell's nil subgroup of the relevant surgery group; $\mathbb Z^-$ denotes the sign-action coefficient module. Its topological realization gives the nonsplittable structures on $\mathbb{RP}^4\#\mathbb{RP}^4$ discussed in the source.
\paragraph{Statement recorded in the supplied source.}
Does the Wall action lift to an action on $\mathcal S_{\mathrm{DIFF}}(X)$ such that
\[
F(A\cdot u)=A\cdot F(u)\qquad
(A\in L_5^s(\mathbb Z[\pi],w),\ u\in\mathcal S_{\mathrm{DIFF}}(X))?
\]
Ask this without stabilizing $X$ by $S^2\times S^2$, in particular for $\pi=\mathbb Z$, $\mathbb Z\times\mathbb Z/n\mathbb Z$ with $n\ge2$, and $\mathbb Z/2\mathbb Z*\mathbb Z/2\mathbb Z$. The source's concrete realization questions are:
\begin{enumerate}
\item[(i)] Can the generator of $L_5(\mathbb Z[\mathbb Z])$ on $S^1\times S^3$ be realized by an exotic smooth $S^1\times S^3$, homeomorphic but not diffeomorphic to the standard product?
\item[(ii)] For $n\ge2$, can its $L_5^s$ action produce a smooth fake $S^1\times L(n,q)$ not diffeomorphic to any product $S^1\times L(n,q')$?
\item[(iii)] Are the nonsplittable topological manifolds in the above nil-group realization on $\mathbb{RP}^4\#\mathbb{RP}^4$ smoothable? Here nonsplittable means not decomposable as a connected sum of two manifolds each homotopy equivalent to $\mathbb{RP}^4$.
\end{enumerate}
The nonorientable third case uses its orientation character in the surgery group. \sref{2898}{2}
\subsection{Short English statement}
Can every topological Wall surgery-group action for a good fundamental group be realized smoothly without stabilization? Include the product, lens-space-product, and nonsplittable real-projective-space cases.
\subsection{Sources}
\begin{description}
\item[\hypertarget{source:2898:1}{S1}] ULAM.AI and the UnsolvedMath Contributors, UnsolvedMath: A Curated Collection of Open Mathematics Problems, record 2898 (KP-4.22). CC BY 4.0. \url{https://huggingface.co/datasets/ulamai/UnsolvedMath}.
\item[\hypertarget{source:2898:2}{S2}] R. İnanç Baykur, Robion C. Kirby and Daniel Ruberman, K3—A New Problem List in Low-Dimensional Topology, Mathematical Surveys and Monographs 295 (2026), author version, Problem 4.22 and Remarks, PDF page 207. \url{https://math.berkeley.edu/sites/default/files/surv-295-ruberman-watermarked-author-pdf.pdf}.
\end{description}
\label{end:2898}
\end{document}
